\specialchapter{历史注记（第七章与第八章）}
% semantic-fragments: legacy-input-seam
\relax{}
(注意 --- 罗马数字指本注末尾的参考文献。)

长度、面积和体积这些概念，在希腊人那里本质上建立在它们对位移的不变性之上：“能重合的东西（εφαρμόξοντα）是相等的”（Eucl. El.，第一卷，“通用概念” 4）；而经典“图形”（多边形、圆锥曲线、多面体、球面等）的面积或体积的所有公式，正是通过对这一原理的巧妙运用而得到的，有时用有限分解的方法，有时用“穷竭法” (*)。用现代的语言可以说，希腊几何学家所做的事情，就是证明“集函数”——可加且在位移下不变——的存在性，只是它们仅对非常特殊类型的集合有定义。积分演算可视为对扩大这类集函数定义域这一需要的回应，而且从 Cavalieri 到 H. Lebesgue，正是这一关切一直处于分析学家研究工作的最前沿；至于位移之下的不变性，它则退居次要地位，因为该性质已成为二重或三重积分中变量替换的一般公式以及正交变换行列式等于 ±1 这一事实的平凡推论。即使在非欧几何中（尽管那里的位移群不同），观点仍然相同：一般地，Riemann 定义了面积或体积的无穷小元素（或维数 \(\geqslant3\) 时它们的类似物），即从一个 \(ds^{2}\) 出发按经典的欧氏公式得到它们，因此它们在使 \(ds^{2}\) 不变的变换下的不变性几乎是同义反复。

直到 1890 年前后，随着积分不变量理论（尤其是 H. Poincaré 与 E. Cartan 的工作）的发展，才出现了群不变测度概念的其他一些不那么直接的推广；H. Poincaré 只考虑了在空间一部分上作用的单参数群，而 E. Cartan 首先关心的是位移群，但作用在与定义它们的那个空间不同的空间中。例如，他由此（除其他结果外）确定了 (II) \(\mathbf{R}^2\) 或 \(\mathbf{R}^3\) 的直线空间上的（在位移群下不变的）不变测度 (*)；此外他还注意到，一般地说，李群上的积分不变量无非是一些特殊的微分不变量，因此可以用 Lie 的方法把它们全部确定出来。然而，在 A. Hurwitz 1897 年的基础性工作 (V) 之前，似乎没有人想到去考虑或使用群本身上的不变测度。为了构造（\(\mathbf{R}^n\) 上）在正交群下不变的多项式，Hurwitz 从这样的评注出发：对线性变换的有限群，只要取变换 \(s \cdot P\) 的平均——其中 \(P\) 是任一多项式，\(s\) 是群的元素——问题便立即解决---这使他想到，对正交群，可用关于一个不变测度的积分来代替平均；他借助 Euler 角给出的参数表示明显地写出了这个测度的表达式，但随即（与 E. Cartan 相互独立地）注意到，Lie 的方法可对每个李群给出不变测度的存在性。也许由于 20 世纪初不变量理论的衰落，Hurwitz 的思想几乎没有得到直接的回响，直到 1924 年以后，随着 I. Schur 与 H. Weyl 把 Frobenius 关于有限群线性表示的经典理论推广到紧群，这些思想才得到发掘。前者限于正交群的情形，并指出 Hurwitz 的方法如何能把特征标的经典正交关系加以推广---H. Weyl 把这一想法与 E. Cartan 关于半单李代数的工作相结合，得到了紧李群不可约表示特征标的明显表达式以及完全可约性定理 (XI \(a\))，然后，通过对“正则表示”概念作一次大胆的推广，得到了著名的 Peter-Weyl 定理，它是有限群理论中正则表示分解为其不可约分量的完美类似 (XI, b))。

早一年，O. Schreier 建立了拓扑群的一般理论，从那时起人们已经清楚，Peter--Weyl 文章中的论证对于任何可在其上定义“不变测度”的拓扑群都原封不动地保持有效。实际上，当时拓扑学与测度论的一般概念仍处于迅速发展之中，哪些拓扑群上可以指望定义不变测度、这种“测度”应当对哪些集合来定义，似乎都尚未明确划定。唯一明显的一点是，不能指望把证明李群上不变测度存在性的无穷小方法推广到一般情形。而另一个思想潮流，从关于 Lebesgue 测度的工作出发，恰好导出了更直接的进攻方法。Hausdorff 在 1914 年证明，不存在定义在 \(R^{3}\) 的所有子集上、在位移下不变且不恒为零的可加集函数，于是自然要问这一结论对 R 与 \(R^{2}\) 是否同样成立：S. Banach 在 1923 年以一种出人意料的方式解决了这个问题，他证明恰恰相反，这样的“测度”确实存在 (I)；他的方法极为巧妙，已经建立在超限归纳构造以及函数被群的元素平移后所得的“平均” \(\frac{1}{n}\sum_{k=1}^{n}f(x+\alpha_{k})\) 的考虑之上 (*)。正是类似的思想使 A. Haar 在 1933 年 (IV) 迈出了决定性的一步，证明了对于开集具有可数基的局部紧群，不变测度的存在性：他受到经典积分演算中用任意小的全等立方体的并置逼近一个体积的方法的引导，借助对角线方法，把不变测度作为一列“近似测度”的极限而得到，这一做法本质上就是我们在 Ch. VII, §1 中所用者。这一发现影响极大，特别是因为它立即使 J. von Neumann 能够就紧群解决 Hilbert 关于用纯拓扑性质（排除任何预先给定的微分结构）刻画李群的著名的“第 5 问题”。然而人们随即察觉，要有效地使用不变测度，不仅需要知道它的存在性，还需要知道它在相差一个常数因子的意义下是唯一的；这一点先由 J. von Neumann 就紧群证明，用的是通过连续函数的（类似于 Banach 的）“平均”定义 Haar 测度的方法 (VII a))；随后 J. von Neumann (VII b)) 与 A. Weil (X) 用不同的方法同时得到了局部紧群情形的唯一性，A. Weil 同时还指明了如何把 Haar 的方法推广到一般的局部紧群。也是 A. Weil（同上）得到了齐性空间上相对不变测度存在的条件，并最终证明，Hausdorff 拓扑群上（具有合理性质的）“测度”的存在性本身就蕴涵该群是局部准紧的。这些工作实质上完成了 Haar 测度的一般理论；此后唯一值得提及的新增内容是拟不变测度的概念，它大约直到 1950 年，随着局部紧群在 Hilbert 空间中的表示理论，才被明确出来。

卷积积的历史更为复杂。早在 19 世纪初人们就注意到，例如，若 \(\mathbf{F}(x,t)\) 是关于 x 与 t 的线性常系数偏微分方程的一个解，则

\[
\int_ {- \infty} ^ {+ \infty} \mathrm{F} (x - s, t) f (s) d s
\]

也是同一方程的解；事实上，早在 1820 年之前，Poisson 等人就已经利用这一想法把热方程的解写成

\[
\int_ {- \infty} ^ {+ \infty} \exp \left(- \frac {(x - s) ^ {2}}{4 t}\right) f (s) d s.\tag{1}
\]

稍后，表达式

\[
\frac {1}{2 \pi} \int_ {- \pi} ^ {+ \pi} \frac {\sin \frac {2 n + 1}{2} (x - t)}{\sin \frac {x - t}{2}} f (t) d t\tag{2}
\]

（它给出 Fourier 级数的部分和）以及 Dirichlet 对该积分当 \(n\) 趋于 \(+\infty\) 时极限的研究，提供了“正则化” \(f \mapsto \rho_n * f\) 在环面 \(\mathbf{T}\) 上的第一个例子（实际上是用一列非正的“核”，这使研究大为复杂化）；这些类似的积分表达式以“奇异积分”之名，从 P. du Bois-Reymond 到 H. Lebesgue，一直是 19 世纪末与 20 世纪初分析学家们偏爱的课题。在 \(\mathbf{R}\) 上，Weierstrass 在其多项式逼近定理的证明中使用了积分 (1)，并在此背景下给出了用正“核”序列 \(\rho_n\) （形如 \(x \mapsto c_n \rho(x/n)\)）作正则化的一般原理。在 \(\mathbf{T}\) 上，用正核作正则化的最著名的例子稍后由 Fejér 给出，从这时起，这一标准程序便成为大多数函数项级数“求和法”的基础。

然而，这些工作由于“核”与被正则化的函数二者所起作用的不对称，几乎没有揭示卷积积的代数性质。我们尤其要感谢 Volterra 强调了这一点。他对二元函数的“复合” \(\mathbf{F}*\mathbf{G}\) 作了一般研究

\[
(\mathbf {F} * \mathbf {G}) (x, y) = \int_ {x} ^ {y} \mathbf {F} (x, t) \mathbf {G} (t, y) d t,
\]

他把它视为两个矩阵之积“由有限到无限的过渡”意义上的推广 (IX)。他很早便单独挑出了 F 与 G 只依赖于 y - x 的情形（由于它在遗传性理论中的解释而称为“闭循环”的情形）；此时 H = F * G 亦然，并且若令 \(\mathrm{F}(x,y)=f(y-x)\) ，\(\mathrm{G}(x,y)=g(y-x)\) ，则

\[
\mathrm{H} (x, y) = h (y - x),
\]

其中

\[
h (t) = \int_ {0} ^ {t} f (t - s) g (s) d s,
\]

于是，当 \(t \geqslant 0\) 时，h 与函数 \(f_{1}, g_{1}\) 的卷积相合，这两个函数当 \(t \geqslant 0\) 时分别等于 f 与 g，而当 t \textless{} 0 时等于 0。

尽管如此，Volterra 所发展的代数形式体系并未揭示出与 \(\mathbf{R}\) 的群结构以及 Fourier 变换的联系。这里不是叙述后者历史的地方；但应当指出，从 Cauchy 起，处理 Fourier 积分的分析学家们主要致力于为各种“反演”公式的有效性寻找越来越弱的条件，而在一定程度上忽视了其代数性质。对 Fourier 本人的工作（以及 Laplace 关于类似积分 \(\int_0^{+\infty}e^{-st}f(t)dt\) 的工作）当然不能这样说；但这些变换本质上是在与线性问题相关的背景下引入的，因此很长时间里无人想到考虑两个 Fourier 变换的乘积也就不足为怪了（三角级数或幂级数的乘积除外，但与离散测度的卷积之间的联系在 19 世纪显然还不可能被察觉）。这一乘积以及 \(\mathbf{R}\) 上卷积的最早提及，大概见于 Tchebychef 的一篇著作 (VIII)，与概率论中的问题有关。事实上，在这门理论中，R 上两个“概率律”（总质量为 1 的正测度）的卷积 \(\mu *\nu\) 无非就是 \(\mu\) 与 \(\nu\) 的“复合”概率律（对相应“随机变量”的加法而言）。诚然，对 Tchebychef 来说，涉及的还只是具有（关于 Lebesgue 测度的）密度的概率律的卷积，因而是函数的卷积；况且它在他的工作中只是附带地出现，在 1920--1930 年之前出现它的那几部罕见著作中情况也是如此。1920 年，P. J. Daniell 在一篇当时很少受人注意的短文 (III) 中定义了 R 上两个任意测度的卷积以及这种测度的 Fourier 变换，并明确指出 Fourier 变换把卷积变成通常的乘积---这一形式体系从 1925 年起被概率学家们密集地使用，尤其是在 P. Lévy 之后。但卷积在群论中的根本重要性直到 1927 年才被 H. Weyl 完全认识到；他注意到，对紧群而言，函数的卷积扮演着有限群代数中乘法的角色，这使他随后得以定义“正则表示”；同时，他在正则化中找到了有限群代数单位元的等价物。剩下的是把所有这些观点加以综合，这由 A. Weil 的书 (X) 完成，它为其后分别构成 I. Gelfand 的赋范代数理论与广义函数卷积的那些推广铺平了道路。

Haar 测度与卷积迅速成为代数化倾向中必不可少的工具，而这一倾向如此强烈地标志着现代分析的特征；我们将在后续各卷中有机会给出它们的众多应用。我们在这些章中处理过的唯一一个应用，涉及局部紧群的闭子群（特别是离散子群）的“变动”。这一理论从 K. Mahler 在数的几何中的一个结果出发，由 C. Chabauty 于 1950 年开创，并且刚刚被 Macbeath 与 Swierczkowski (VII) 大大地发展与深化，我们在此复述了他们的主要结果。

